% Options for packages loaded elsewhere
\PassOptionsToPackage{unicode}{hyperref}
\PassOptionsToPackage{hyphens}{url}
\documentclass[
]{article}
\usepackage{xcolor}
\usepackage{amsmath,amssymb}
\setcounter{secnumdepth}{-\maxdimen} % remove section numbering
\usepackage{iftex}
\ifPDFTeX
  \usepackage[T1]{fontenc}
  \usepackage[utf8]{inputenc}
  \usepackage{textcomp} % provide euro and other symbols
\else % if luatex or xetex
  \usepackage{unicode-math} % this also loads fontspec
  \defaultfontfeatures{Scale=MatchLowercase}
  \defaultfontfeatures[\rmfamily]{Ligatures=TeX,Scale=1}
\fi
\usepackage{lmodern}
\ifPDFTeX\else
  % xetex/luatex font selection
\fi
% Use upquote if available, for straight quotes in verbatim environments
\IfFileExists{upquote.sty}{\usepackage{upquote}}{}
\IfFileExists{microtype.sty}{% use microtype if available
  \usepackage[]{microtype}
  \UseMicrotypeSet[protrusion]{basicmath} % disable protrusion for tt fonts
}{}
\makeatletter
\@ifundefined{KOMAClassName}{% if non-KOMA class
  \IfFileExists{parskip.sty}{%
    \usepackage{parskip}
  }{% else
    \setlength{\parindent}{0pt}
    \setlength{\parskip}{6pt plus 2pt minus 1pt}}
}{% if KOMA class
  \KOMAoptions{parskip=half}}
\makeatother
\usepackage{longtable,booktabs,array}
\usepackage{caption}
\captionsetup[table]{skip=6pt}
\newcounter{none} % for unnumbered tables
\usepackage{calc} % for calculating minipage widths
% Correct order of tables after \paragraph or \subparagraph
\usepackage{etoolbox}
\makeatletter
\patchcmd\longtable{\par}{\if@noskipsec\mbox{}\fi\par}{}{}
\makeatother
% Allow footnotes in longtable head/foot
\IfFileExists{footnotehyper.sty}{\usepackage{footnotehyper}}{\usepackage{footnote}}
\makesavenoteenv{longtable}
\usepackage{graphicx}
\makeatletter
\newsavebox\pandoc@box
\newcommand*\pandocbounded[1]{% scales image to fit in text height/width
  \sbox\pandoc@box{#1}%
  \Gscale@div\@tempa{\textheight}{\dimexpr\ht\pandoc@box+\dp\pandoc@box\relax}%
  \Gscale@div\@tempb{\linewidth}{\wd\pandoc@box}%
  \ifdim\@tempb\p@<\@tempa\p@\let\@tempa\@tempb\fi% select the smaller of both
  \ifdim\@tempa\p@<\p@\scalebox{\@tempa}{\usebox\pandoc@box}%
  \else\usebox{\pandoc@box}%
  \fi%
}
% Set default figure placement to htbp
\def\fps@figure{htbp}
\makeatother
\ifLuaTeX
  \usepackage{luacolor}
  \usepackage[soul]{lua-ul}
\else
  \usepackage{soul}
\fi
\setlength{\emergencystretch}{3em} % prevent overfull lines
\providecommand{\tightlist}{%
  \setlength{\itemsep}{0pt}\setlength{\parskip}{0pt}}
\ifPDFTeX
  \TeXXeTstate=1
  \newcommand{\RL}[1]{\beginR #1\endR}
  \newcommand{\LR}[1]{\beginL #1\endL}
  \newenvironment{RTL}{\beginR}{\endR}
  \newenvironment{LTR}{\beginL}{\endL}
\fi
\ifluatex
  \newcommand{\RL}[1]{\bgroup\textdir TRT#1\egroup}
  \newcommand{\LR}[1]{\bgroup\textdir TLT#1\egroup}
  \newenvironment{RTL}{\textdir TRT\pardir TRT\bodydir TRT}{}
  \newenvironment{LTR}{\textdir TLT\pardir TLT\bodydir TLT}{}
\fi
\usepackage{bookmark}
\IfFileExists{xurl.sty}{\usepackage{xurl}}{} % add URL line breaks if available
\urlstyle{same}
% fallback for those not using the hyperref driver hyperxmp:
\makeatletter
\@ifundefined{xmpquote}{\newcommand{\xmpquote}[1]{#1}}{}
\makeatother
\hypersetup{
  hidelinks,
  pdfcreator={LaTeX via pandoc}}

\author{}
\date{}

\begin{document}

\includegraphics[width=2.62865in,height=1.57545in]{media/image1.png}\LR{}

\LR{}

\LR{}

Thesis Progress Report

Predicting Clinical Outcomes and Assessing the Effects of Diabetes Care
Guidelines Compliance in Hospitals Using Time-Dependent Machine Learning
and Deep Learning Methods

January 2026

\textbf{Submitted by:}

Shahar Oded

Faculty of Computer and Information Science

Ben-Gurion University of the Negev

E-Mail:
\href{mailto:odedshah@post.bgu.ac.il}{\ul{odedshah@post.bgu.ac.il}}

\textbf{Supervised by:}

Prof. Yuval Shahar

Faculty of Computer and Information Science

Ben-Gurion University of the Negev

E-Mail: \href{mailto:yshahar@bgu.ac.il}{\ul{yshahar@bgu.ac.il}}

\ul{Dataset:}

Significant progress has been made in the end-to-end data engineering
pipeline, transforming raw administrative and clinical data into a
research-ready, longitudinal format. This phase focused on the
\emph{Kineret} dataset, comprising data from six government hospitals,
and involved the following key achievements:

\begin{itemize}
\item
  \textbf{ETL and Indexing:} A robust Extract, Transform, Load (ETL)
  process was developed to map the source data into a standardized
  longitudinal format. Since the raw data lacked international coding
  standards, a comprehensive \textbf{concept index} of \textasciitilde2K
  rows was constructed to map all relevant local codes to unified
  clinical concepts, including unit normalization. This enabled the
  ingestion of heterogeneous data sources into a single, cohesive
  structure.
\item
  \textbf{Knowledge Base (KB) and Discretization:} A machine-readable KB
  was defined to support Temporal Abstraction. To ensure clinical
  validity, \textbf{discretization rules} were established using a
  hybrid approach: data-driven K-means clustering, KDE density plots
  were used to identify natural data thresholds, which were then refined
  by real world medical knowledge. The KB currently encompasses
  definitions for \textbf{raw concepts, states, events, trends, and
  contexts}. The knowledge base follows a custom scheme I developed to
  support my treatment guidelines requirements, and data manipulation
  goals, turning the data to pluggable input to Machine Learning and
  Deep Learning architectures.
\item
  \textbf{Treatment Patterns and QA Scoring:} Beyond basic abstractions,
  the KB now includes \textbf{treatment patterns} that define "correct"
  therapeutic actions according to the clinical guidelines.
  \textbf{Fuzzy logic scores} were implemented to quantify adherence to
  these patterns. These patterns are outputted as a part of the
  longitudinal dataset, accompanied by a compliance scores table that
  includes the fuzzy scores. \emph{Note: As the clinical protocol
  evolves, these specific pattern definitions and scoring mechanisms
  remain adaptable to future updates.}
\item
  \textbf{Mediator Implementation and Context Vector:} A version of the
  \textbf{Temporal Mediator} system was successfully implemented to
  apply these KB rules to the raw data. This system outputs an
  informative, interval-based longitudinal dataset. Additionally, a
  parallel ETL process was built to generate a static \textbf{context
  vector} for each hospitalization, capturing demographic information,
  essential background data, the Charlson Comorbidity Index (CCI), and
  relevant home medications.
\item
  \textbf{Data Reduction and Optimization:} The efficacy of this
  pipeline is demonstrated by the massive reduction in data noise and
  dimensionality. The process successfully filtered and abstracted
  \textbf{hundreds of millions of raw records} down to a streamlined,
  informative table of approximately \textbf{10 million interval
  records}, which includes the calculated Guideline (GL) adherence
  scores for every admission.
\end{itemize}

\ul{Architectures:\\
}Parallel to the data processing efforts, the computational
architectures for both the frequent-pattern-based approach and the deep
sequential modeling approach have been fully implemented and are ready
for the experimental phase.

\textbf{1. Vector-Space Modeling (VSM): The KarmaLego Engine}

For the baseline comparison, a complete, end-to-end Python
implementation of the \textbf{KarmaLego} algorithm has been developed to
mine frequent Time Interval Relation Patterns (TIRPs). This engine
transforms the longitudinal interval data into fixed-length feature
vectors suitable for standard machine learning classifiers. Key
technical achievements include:

\begin{itemize}
\item
  \textbf{Pipeline Architecture:} The system manages the full lifecycle
  from data ingestion to pattern application. It executes the "Karma"
  phase (temporal relation encoding) and the "Lego" phase (tree-based
  pattern extension) to discover high-order temporal patterns. After
  patterns are globally detected, the engine applies them on the
  patient's dataset (see ``Feature Generation'').
\item
  \textbf{Performance Optimization:} To handle the scale of the EMR
  dataset, the implementation incorporates \textbf{Apriori pruning} and
  \textbf{memoized temporal relation transitivity}. Crucially, it
  implements \textbf{Consistent Subset of Active Candidates (CSAC)}, a
  rigorous pruning method that maintains exact parent embedding indices
  to avoid searching impossible extensions, ensuring scalability on
  dense timelines. Target oriented parallelism was also added.
\item
  \textbf{Feature Generation:} The pipeline is capable of applying
  discovered patterns back to the patient cohort to generate
  \textbf{Temporal Probabilistic Profiles (TPF)}. It supports multiple
  output modes, including pattern counts, cohort-normalized
  distributions, and duration-based features, effectively creating the
  input vectors required for classical ML experiments.
\end{itemize}

\textbf{2. Deep Sequential Modeling: The Transform-EMR Architecture}

A novel, two-phase deep learning pipeline, titled
\textbf{Transform-EMR}, has been finalized to model the complex
dependencies in the EMR data (Mediator's output). This architecture is
designed to predict patient event trajectories by explicitly modeling
time, hierarchy, and patient context.

\begin{itemize}
\item
  \textbf{Phase 1: Representation Learning (EMREmbedding):} A dedicated
  embedding module was built to isolate the core structure of patient
  timelines.

  \begin{itemize}
  \item
    \textbf{Hierarchical Tokenization:} Events are decomposed into four
    hierarchical components (Raw Concept -\textgreater{} Concept
    -\textgreater{} Value -\textgreater{} Position) to enforce
    similarity across related medical tokens.
  \item
    \textbf{Temporal \& Contextual Fusion:} Time is encoded using a
    \textbf{Time2Vec} layer (learning periodic and linear trends) and
    concatenated with event tokens via Early Fusion. Static patient
    context (e.g., demographics) is projected and added to the sequence
    as a global bias.
  \item
    \textbf{Pre-training Objectives:} The embedder is trained using a
    multi-objective loss function combining Next-Token Prediction (BCE),
    Time Prediction (MSE), and Masked Language Modeling (MLM).
  \end{itemize}
\item
  \textbf{Phase 2: Sequence Modeling (Transformer Decoder):} The
  generative model is built upon a \textbf{GPT-style decoder-only
  architecture} tailored for clinical sequences.

  \begin{itemize}
  \item
    \textbf{Context Conditioning:} Standard normalization layers were
    replaced with \textbf{AdaLN-Zero (Adaptive Layer Norm)} blocks,
    allowing the static patient context vector to dynamically modulate
    the generation process at every transformer block.
  \item
    \textbf{Constrained Training \& Inference:} The training loop
    utilizes a weighted focal loss combined with \textbf{structural
    penalties} (e.g., enforcing logical interval starts/ends and meal
    ordering). An inference engine was developed to support step-by-step
    generation with dynamic masking, ensuring the model cannot predict
    "illegal" tokens that violate the temporal logic of the generated
    stream.
  \end{itemize}
\end{itemize}

\ul{Analysis:}

To validate the clinical relevance of the generated data and adherence
scores, we analyzed the abstracted output from the Mediator. This
analysis, \textbf{submitted to the \emph{Israeli Endocrinology
Conference 2026}}, utilized the \emph{Kineret} dataset to examine the
link between guideline adherence and hospital-acquired complications.

\begin{itemize}
\item
  \textbf{Methodology:} We compared adherence scores - calculated via
  the framework described above - between the pre- and
  post-implementation periods of the national guidelines.
\item
  \textbf{Key Findings:} The analysis revealed that higher adherence to
  the protocols was associated with lower rates of in-hospital
  mortality. Specific improvements were observed in the discontinuation
  of high-risk medications and insulin de-escalation following
  hypoglycemia.
\item
  \textbf{Predictive Validation:} Using bootstrapped elastic-net
  multivariate models, we found that guideline adherence scores
  accounted for up to \textbf{40\% of the predictive weight} for
  hospital-acquired complications. This strongly validates the
  hypothesis that the QA scores generated by our pipeline are critical,
  high-signal features for the subsequent deep learning prediction
  tasks.
\end{itemize}

\ul{Future Timeline:\\
}Building on the completed architectures and the preliminary clinical
analysis, the remaining research phase will focus on rigorous
experimentation and benchmarking:

{\def\LTcaptype{none} % do not increment counter
\begin{longtable}[]{@{}
  >{\raggedright\arraybackslash}p{(\linewidth - 2\tabcolsep) * \real{0.1811}}
  >{\raggedright\arraybackslash}p{(\linewidth - 2\tabcolsep) * \real{0.8189}}@{}}
\toprule\noalign{}
\endhead
\bottomrule\noalign{}
\endlastfoot
Period & Milestone \& Objectives \\
Jan -- Mid-Feb 2026 & Execute core experiments comparing Transform-EMR
vs. TPF baselines and possibly additional SOTA models like STraTS /
BEHRT/ TransformEHR.

Finalize prediction quality metrics for the "AIME" conference or
workshop deadline.

Submit preliminary algorithmic results. \\
Feb -- March 2026 & ablation studies (Impact of QA scores/Time2Vec).

Explainability Analysis: Perform SHAP analysis on the finalized model.

Draft the "Results" and "Discussion" chapters of the thesis. \\
April -- July 2026 & \textbf{Complete full thesis manuscript and
presentation.}

Refine the thesis work into camera-ready papers for major conference
submissions.

Submit final manuscripts to top-tier venues (e.g., IEEE-BIBM 2026, AAAI
2027) dividing the work to architectural novelty and to medical
implementation. \\
August 2026 & Formal submission of the Master\textquotesingle s Thesis +
defense. \\
\end{longtable}
}

\LR{}

\end{document}
